\chapter{Uniform bounds for inserting twist powers}
\label{app:uniform-filling-bound}

The finite formula for a diagram can be organized in two steps. First,
take the prescribed coefficient in its shared face variables. This gives
a polynomial \(S_N\) in one spectral variable for each crossing region.
Then apply a linear functional in each spectral variable to insert the
chosen signed twist power.

This appendix bounds the second step. Inserting \(n\) crossings costs at
most \(1+|n|/2\) in the supremum norm on the real spectral interval,
independently of the quantum level \(N\). There is a stronger estimate
for the two one-sided pieces of the local polynomial: their values
vanish at all spectral nodes, and the cost is then at most \((N-1)/N\),
independently of \(n\) as well. The derivatives at repeated nodes remain
in the calculation.

The estimate applies after the face coefficient has been taken. All
crossing regions therefore retain their common face constraints. To
obtain a volume theorem from it, one must still bound this shared
polynomial and prove that the geometric contribution survives.

Here is the precise statement. Let \(N\ge2\), \(d=N-1\),
\(q=e^{\pi i/N}\), and let \(\mathcal L_N^{n}\) be the full-Hermite
functional for the actual crossing power \(R^n\), with
\(n\in\mathbb Z\), defined below. For a polynomial \(H\) of degree at
most \(d\), \begin{equation}\label{uf:eq-1}
|\mathcal L_N^{n}H|
       \le \left(1+\frac{|n|}{2}\right)
                         \|H\|_{[-2,2]}.
\end{equation} If \(H\) vanishes at every root of \(\chi_N(z)=2T_N(z/2)+2\),
including the simple root \(-2\) when \(N\) is odd, the stronger bound
is \begin{equation}\label{uf:eq-2}
|\mathcal L_N^{n}H|\le \frac dN\|H\|_{[-2,2]}.
\end{equation} Vanishing means values only; the derivatives at double roots are
retained.

For \(r\) crossing regions, let \(S(\sigma_1,\ldots,\sigma_r)\) have
degree at most \(d\) in each variable. Suppose that for every \(e\) in a
set \(J\), \(S\) vanishes whenever \(\sigma_e\) is a spectral node, with
its other variables arbitrary. Then \begin{equation}\label{uf:eq-3}
\left|\left(\bigotimes_{e=1}^r\mathcal L_N^{n_e}\right)S\right|
 \le \left(\frac dN\right)^{|J|}
       \prod_{e\notin J}\left(1+\frac{|n_e|}{2}\right)
                          \|S\|_{[-2,2]^r}.
\end{equation} In particular this applies to each term containing a specified lower
or upper one-sided factor, with \(J\) its one-sided crossing regions.
The case \(J=\varnothing\) is the full-source bound.

There is also a bound for the whole complementary source without
expanding it into sectors. Let
\(\mathcal D(\sigma_1,\ldots,\sigma_r)\) have degree at most
\(d=N-1\) in each variable. If \(\mathcal D\) vanishes on the Cartesian
product of all spectral node sets, then \begin{equation}\label{uf:eq-4}
\left|\left(\bigotimes_{e=1}^r\mathcal L_N^{n_e}\right)\mathcal D\right|
 \le \left\{\prod_e\left(1+\frac{|n_e|}{2}\right)
       -\prod_e\left(\frac1N+\frac{|n_e|}{2}\right)\right\}
                         \|\mathcal D\|_{[-2,2]^r}.
\end{equation} For the actual numerator one may take
\(\mathcal D=S_{\mathrm{full}}-S_{\mathrm C}\), where
\(S_{\mathrm C}\) retains the central piece at every crossing region
(defined in Section~\ref{mca:sec-1}). The coefficient in
braces has degree at most \(r-1\) in the quantities \(|n_e|\): the
highest product cancels exactly. None of \eqref{uf:eq-1}--\eqref{uf:eq-4} assumes positivity of
spectral coefficients, asymptotic localization, or dominance of a
Fourier charge.

\section{The exact source and its
normalization}\label{uf:sec-the-exact-source-and-its-normalization}

Use the original pair Casimir \(\Omega\), not an operator reconstructed
from its spectral values. The direct matrix calculation in
Section~\ref{uf:local-fourier} proves \(\chi_N(\Omega)=0\), and
polynomial division gives the degree-\(d\) resolvent numerator \begin{equation}\label{uf:eq-5}
V_N(z;\Omega)=\chi_N(z)(z-\Omega)^{-1}.
\end{equation} The right side means the resulting matrix polynomial, including at
every spectral node. It is not the determinant adjugate of the full
\(N^2\)-dimensional matrix.

The
\href{https://sashakolpakov.github.io/murakamis_dream/notes/actual_homogeneous_casimir_face_source.html}{original
homogeneous face source} associates to \eqref{uf:eq-5} a four-face polynomial
\(H_V(x,b,c,y;z)\), with literal input/output middle faces \(b,c\). Fix
a connected original diagram at \(N\ge3\), with its actual cups, caps,
unrestricted marked cap, protected face targets, and every ordered box.
The applications include connected flat plats and the connected matched
simple-sphere Morse diagrams of
\href{https://sashakolpakov.github.io/murakamis_dream/notes/research_matched_sphere_morse_source.html}{the
general Morse source theorem}, §§1,5--6, and
\href{https://sashakolpakov.github.io/murakamis_dream/notes/actual_balanced_terminal_fan_volume.html}{the
terminal-fan source proof}, §§5--7. For the latter, the actual cyclic
orders and crossing-disk frames are retained, and the marked main edge
determines the two protected faces. The exact coefficient identity is
the only input here: no complete-radius balance or geometric certificate
is required for this finite norm estimate. The shared spectral
polynomial is \begin{equation}\label{uf:eq-6}
S_N(\boldsymbol\sigma)=
 [\textstyle\prod_{f\ \mathrm{ordinary}}L_f^d]
       \prod_{e=1}^{r}H_V(L_{x_e},L_{b_e},L_{c_e},L_{y_e};\sigma_e),
\end{equation} with the two protected face exponents zero. Every variable
\(\sigma_e\) has degree at most \(d\). The original marked source after
inserting the signed powers is exactly \begin{equation}\label{uf:eq-7}
Z_N(\boldsymbol n)=
        \left(\bigotimes_{e=1}^r\mathcal L_N^{n_e}\right)S_N.
\end{equation} The face-grid factors, cup/cap factors and the original \(N^{-1}\)
marked scalar have already canceled as in those source theorems; there
is no additional \(N^r\) in \eqref{uf:eq-7}. Thus \(Z_N\) is the original reduced
ribbon contraction. Conversion to the prescribed zero-framed enhanced
Kashaev convention uses the recorded framing/component character, of
modulus one at complete color; for a knot the enhanced component sign is
one. In the specified full one-cut modified convention, retain the
additional factor \(i^{N+3}/N\). Original exterior Cartan and framing
factors remain in their prescribed locations throughout. These are the
physical normalization choices; no new quantum dimension or color
average is introduced by the norm theorem.

The polynomial degree and the identities \eqref{uf:eq-5}--\eqref{uf:eq-7}, rather than a
hypothesis about arbitrary polynomial solutions of a recurrence, are the
source assumptions used here. The scalar inequalities \eqref{uf:eq-1}--\eqref{uf:eq-4}
themselves hold for every polynomial of the indicated degree and every
\(N\ge2\).

\section{The local source identity from a finite Fourier transform}
\label{uf:local-fourier}

We give the local matrix calculation, including its normalization.
It explains both the annihilating polynomial and the homogeneous face
polynomial used above. No decomposition into simple modules is needed.

Put \(Q=q^2\). In the balanced basis \(e_0,\ldots,e_d\) of the
complete-color module, the original matrices are
\[
 K e_i=q^{d-2i}e_i,\qquad
 F e_i=(1-Q^{i+1})e_{i+1},\qquad
 \mathcal E e_i=(Q^{-i}-1)e_{i-1}.
\]
Here \(\mathcal E=(q-q^{-1})^2E\), and a term outside
\(0,\ldots,d\) is zero. These formulas follow by replacing the
usual basis vector \(v_i\) by \(v_i/(Q;Q)_i\); all these
denominators are nonzero because \(i<N\). The pair matrices are
\[
 \begin{split}
 K_{12}&=K\otimes K,\\
 F_{12}&=F\otimes I+K^{-1}\otimes F,\\
 \mathcal E_{12}&=\mathcal E\otimes K+I\otimes\mathcal E,\\
 \Omega&=F_{12}\mathcal E_{12}
                     +qK_{12}+q^{-1}K_{12}^{-1}.
 \end{split}
\]
On an input \(e_u\otimes e_v\), set \(U=Q^u,V=Q^v\).
Multiplying these displayed matrices gives the diagonal coefficient
\(D\), and the two coefficients \(A_+,A_-\) that change the first
index by \(+1,-1\):
\begin{equation}\label{uf:casimir-three-entries}
\begin{split}
 D(U,V)&=2qU+2q^{-1}V^{-1}-(q+q^{-1})U/V,\\
 A_+(U,V)&=(V^{-1}-1)(1-QU),\\
 A_-(U,V)&=(Q^{-1}V^{-1}-1)(1-U).
\end{split}
\end{equation}
The second index changes by the opposite amount. The positive shift
vanishes when \(u=d\) or \(v=0\), and the negative shift vanishes
when \(u=0\) or \(v=d\). Thus these formulas retain the true
integer total and all physical endpoints.

For independent outer face variables \(x,y\), define an operator on
polynomials in \(z\) by
\begin{equation}\label{uf:face-difference-operator}
\begin{split}
 (Tg)(z)={}&-q^{-1}
   \left(1+\frac{q(x+y)}z+\frac{q^2xy}{z^2}\right)g(Q^{-1}z)\\
 &-q\left(1+\frac{q^{-1}(x+y)}z
                     +\frac{q^{-2}xy}{z^2}\right)g(Qz)\\
 &+\left(\frac{2(x+y)}z
                 +\frac{(q+q^{-1})xy}{z^2}\right)g(z).
\end{split}
\end{equation}
Its denominators cancel on polynomials. Indeed,
\[
 Tz^j=A_jz^j+B_j(x+y)z^{j-1}+C_jxyz^{j-2},
\]
where
\[
 \begin{split}
 A_j&=-(qQ^j+q^{-1}Q^{-j}),\\
 B_j&=2-Q^j-Q^{-j},\\
 C_j&=q+q^{-1}-qQ^{-j}-q^{-1}Q^j.
 \end{split}
\]
The only potentially negative powers have coefficients
\(B_0=C_0=C_1=0\). In particular \(T\) preserves the
\(N\)-dimensional space of polynomials of degree at most \(d\),
and it preserves total degree when \(x,y,z\) all have degree one.

\begin{proposition}[The actual Casimir and the face operator]
\label{prop:local-face-fourier}
For \(s=0,\ldots,d\), let \(\Omega_s\) be the restriction of
\(\Omega\) to the ordered basis
\[
 e_u\otimes e_{[s-u]_N},\qquad 0\le u<N,
\]
where \([a]_N\) is the representative of \(a\) in
\(\{0,\ldots,N-1\}\). Choose \(y\ne0\), set
\(x=Q^{s+1}y\), and put
\[
 b_u=-q^{-1}xQ^{-u},\qquad
 \mathsf E_{uj}=b_u^j,\qquad
 G_{uu}=q^{du}.
\]
The matrix \(\mathsf E\) evaluates a polynomial at the \(N\)
distinct points \(b_u\). With \(T\) written in the monomial basis,
\begin{equation}\label{uf:face-fourier-conjugacy}
 \mathsf E T\mathsf E^{-1}=-G^{-1}\Omega_s^{\mathsf T}G.
\end{equation}
Consequently the full original pair matrix satisfies
\(\chi_N(\Omega)=0\), including its nilpotent parts.
\end{proposition}

\begin{proof}
At \(z=b_u\), the two physical clocks are
\[
 -q^{-1}x/b_u=Q^u,\qquad
 -q^{-1}b_u/y=Q^{s-u}.
\]
The coefficients in \eqref{uf:face-difference-operator} are therefore
\(-D\), \(q^{-1}A_+\), and \(qA_-\). Since
\(q^d=-q^{-1}\), the positive coefficient is
\(-q^dA_+\), and the negative one is \(-q^{-d}A_-\).
Evaluation at \(Q^{-1}b_u\) increments the first index, while
evaluation at \(Qb_u\) decrements it. The endpoint zeros in
\eqref{uf:casimir-three-entries} kill every attempted wrap.
Thus the coefficient from row \(u\) to column \(r\) is
\(-q^{d(r-u)}(\Omega_s)_{ru}\), proving
\eqref{uf:face-fourier-conjugacy}. The transpose occurs because
evaluation acts on functions of the input index, whereas the physical
matrix sends input columns to output rows.

The matrix of \(-T\) is triangular, with diagonal
\(t_j+t_j^{-1}\), where \(t_j=qQ^j\). Its characteristic
polynomial is
\[
 \prod_{j=0}^{N-1}(X-t_j-t_j^{-1})=2T_N(X/2)+2.
\]
To check the identity with multiplicities, substitute \(X=w+w^{-1}\)
and multiply by \(w^N\). The left side becomes
\[
 \prod_{j=0}^{N-1}(w-t_j)(w-t_j^{-1})=(w^N+1)^2,
\]
because both lists are exactly the roots of \(t^N=-1\).
The right side gives the same Laurent polynomial. Substitution
\(X=w+w^{-1}\) is injective on polynomials, so the identity follows.
Cayley--Hamilton and \eqref{uf:face-fourier-conjugacy} give
\(\chi_N(\Omega_s)=0\). The \(N\) residue spaces form a direct sum
of the full pair space, proving the assertion for \(\Omega\).
This proves a matrix identity, not only an identity on its eigenvalues.
\end{proof}

The same transform gives the complete local polynomial adapter.
For any polynomial \(f\), write
\[
 f(-T)z^j=\sum_{\alpha,\beta\ge0,\ \alpha+\beta\le j}
  M_f(j;\alpha,\beta)x^\alpha y^\beta z^{j-\alpha-\beta}
\]
and define
\begin{equation}\label{uf:homogeneous-local-adapter}
 H_f(x,b,c,y)=\sum_{j=0}^{d}\,
 \sum_{\alpha+\beta\le j}
 M_f(j;\alpha,\beta)
 x^\alpha b^{j-\alpha-\beta}c^{d-j}y^\beta.
\end{equation}
Its exponents are nonnegative and its total degree is \(d\),
regardless of the degree of \(f\). Fourier inversion is particularly
simple:
\[
 (\mathsf E^{-1})_{jr}=\frac1N b_r^{-j}.
\]
For a polynomial \(f\), \eqref{uf:face-fourier-conjugacy} gives
\[
 \mathsf E f(-T)\mathsf E^{-1}
                  =G^{-1}f(\Omega_s)^{\mathsf T}G.
\]
Use the displayed inverse of \(\mathsf E\). For \(b=b_u,c=b_r\),
the result is
\begin{equation}\label{uf:physical-local-adapter}
 (f(\Omega_s))_{ru}
   =\frac1N q^{-d(r-u)}c^{-d}H_f(x,b,c,y).
\end{equation}
This proves the factor \(1/N\), the direction of the shift and the
full polynomial identity directly from the original pair matrices.
If the input and output have different integer totals, the left
side is zero: the modular block is the direct sum of its actual-total
blocks. Hence the right side also vanishes at such a grid point.

At even levels \(q^{d(p+N)}=-q^{dp}\). The exponent \(r-u\)
in \eqref{uf:physical-local-adapter} is therefore the canonical
integer displacement, not an arbitrarily chosen residue. The
endpoint calculation above is what permits the Fourier transform;
it does not assume that this phase is periodic. With the coherent
face-logarithm lifts, \(c=Q^{-(r-u)}b\) converts the prefactor to
\(N^{-1}(bc)^{-d/2}\), exactly as in the global source calculation.

This argument holds at every fixed root with \(N\ge2\), including
composite levels. It proves the annihilator and the adapter for
polynomials of the actual Casimir. The polynomial representing a
crossing power, its repeated-node derivatives, and the comparison
with the specified knot invariant are separate inputs. No derivative
with respect to a generic deformation of \(q\) is taken here.

\section{Why the polynomial reduction is well defined}
\label{uf:sec-canonical-reduction}

Before using the particular spectral roots, we record the algebraic
fact behind reduction at repeated roots. It explains which parts of the
calculation depend only on the polynomial and which still require the
original crossing operator.

Let \(R\) be a commutative ring with identity, and let
\(r_1,\ldots,r_m\in R\), with \(m\ge1\). Repetitions are allowed.
Put \(p(X)=\prod_{j=1}^m(X-r_j)\). Given \(f\in R[X]\), define
\(f_0=f\) and successively define \(f_j\) by
\[
 f_{j-1}(X)=f_{j-1}(r_j)+(X-r_j)f_j(X),\qquad 1\le j\le m.
\]
Each \(f_j\) is obtained by synthetic division. This is a polynomial
operation: it uses no inverse in \(R\), even if two roots coincide.
Define
\[
 \mathcal R_p f(X)=
 \sum_{j=1}^m f_{j-1}(r_j)\prod_{i<j}(X-r_i),
 \qquad
 \ell_p(f)=[X^{m-1}]\mathcal R_p f.
\]
Here \([X^k]\) means the coefficient of \(X^k\), and an empty
product is one.

\begin{proposition}[Canonical polynomial remainder]
\label{prop:canonical-polynomial-remainder}
The polynomial \(\mathcal R_p f\) is the unique polynomial of degree
less than \(m\) for which \(f-\mathcal R_p f\) is a multiple of
\(p\). Consequently it is independent of the ordering of the root
list, including repeated entries. Both \(\mathcal R_p\) and
\(\ell_p\) are \(R\)-linear, and adding a multiple of \(p\) to
their input leaves their values unchanged. Moreover,
\[
 \ell_p(f)=f_{m-1}(r_m).
\]
Among additive functionals on \(R[X]\), \(\ell_p\) is uniquely
specified by the two requirements
\[
 F(pq)=0\quad\text{for every }q\in R[X],\qquad
 F(r)=[X^{m-1}]r\quad\text{whenever }\deg r<m.
\]
\end{proposition}

\begin{proof}
Substituting each synthetic-division identity into the preceding one
gives \(f=p f_m+\mathcal R_p f\). Every summand in the displayed
remainder has degree at most \(m-1\). If two such remainders existed,
their difference would have degree less than \(m\) and would be a
multiple of the monic polynomial \(p\). A nonzero multiple \(pq\)
has degree \(m+\deg q\): its leading coefficient is the leading
coefficient of \(q\), because that of \(p\) is one. This argument
works even when \(R\) has zero divisors. The difference must therefore
vanish.

Uniqueness proves independence of root order, linearity and invariance
under adding a multiple of \(p\). In the Newton expression for the
remainder, only the last summand can contribute to \(X^{m-1}\);
its product of linear factors is monic. This proves the formula for
\(\ell_p\). Finally, any additive \(F\) with the two stated
properties satisfies
\[
 F(f)=F(pf_m)+F(\mathcal R_p f)
      =[X^{m-1}]\mathcal R_p f=\ell_p(f).
\]
\end{proof}

For two equal roots this construction gives
\[
 \mathcal R_{(X-r)^2}f(X)=f(r)+f'(r)(X-r),
 \qquad \ell_{(X-r)^2}(f)=f'(r).
\]
The derivative is the formal polynomial derivative. Thus repeated
roots retain derivative information; replacing them by a single
evaluation changes the operation.

\begin{proposition}[Products and repeated-root derivatives]
\label{prop:canonical-remainder-products}
For any $f,g\in R[X]$,
\[
 \mathcal R_p(fg)=
 \mathcal R_p\bigl((\mathcal R_p f)(\mathcal R_p g)\bigr),
 \qquad
 \ell_p(fg)=\ell_p\bigl((\mathcal R_p f)(\mathcal R_p g)\bigr).
\]
Consequently a finite product may be reduced after each multiplication,
with the same final remainder and coefficient functional. This reduced
multiplication is associative. If $r$ occurs at least twice among the
roots of $p$, the final remainder retains both $(fg)(r)$ and
\[
 (fg)'(r)=f'(r)g(r)+f(r)g'(r).
\]
For the modulus $(X-r)^2$ its full expression is
\[
 \mathcal R_{(X-r)^2}(fg)
 =f(r)g(r)+\bigl(f'(r)g(r)+f(r)g'(r)\bigr)(X-r).
\]
No degree bound on the unreduced factors or their product is required.
\end{proposition}

\begin{proof}
Write $f=\mathcal R_p f+pa$ and $g=\mathcal R_p g+pb$ using
synthetic division. Their product differs from
$(\mathcal R_p f)(\mathcal R_p g)$ by
$p\{a\mathcal R_p g+b\mathcal R_p f+pab\}$. Uniqueness of the
remainder proves the first equality; taking its top coefficient proves
the second. Repeating the equality proves the finite-product statement.
Either parenthesization of a product of three factors reduces the same
polynomial $fgh$, which proves associativity.

When $(X-r)^2$ divides $p$, the difference between any polynomial and its
remainder has both value and formal derivative zero at $r$. The formal
product rule gives the displayed derivative. The explicit degree-one
remainder then follows from the double-root formula above.
\end{proof}

The coefficient functional is linear, but it is generally not
multiplicative: $\ell_p(fg)$ need not equal $\ell_p(f)\ell_p(g)$.
The product must be formed and reduced before its coefficient is taken.
At a repeated spectral root the product rule also retains the first
derivative of either factor; a list of values alone cannot replace it.
This is the algebra needed when a product initially exceeds the
canonical degree in Appendix~\ref{app:directed-all-orders}.

These polynomial statements are checked in the Lean files described
in Chapter~\ref{ch:reproduction}. The implementation uses coefficient
lists and also proves invariance under arbitrary trailing zeros.
The characterization of \(\ell_p\) states sufficient conditions
for identifying another functional with it; those conditions must
still be proved for that functional. In this appendix the application
requires the actual factorization of \(\chi_N\), the original
Casimir relation, and the crossing values and derivatives. The next
section supplies the root data and the corresponding explicit formula.

There is also a direct algebraic relation between this coefficient
functional and the resolvent numerator. It does not require a
diagonalization of the operator.

\begin{proposition}[Reproduction by the polynomial resolvent]
\label{prop:resolvent-reproduction}
Let \(p(X)\) be a split monic polynomial of degree \(m\ge1\)
over a commutative algebra \(A\), and let \(w\in A\) satisfy
\(p(w)=0\). Define \(V_w\) by polynomial synthetic division,
\[
 V_w(X)=\frac{p(X)-p(w)}{X-w}.
\]
The fraction denotes a polynomial, with no inverse assumed in \(A\).
Let \(\ell_p\) be the coefficient of \(X^{m-1}\) in the
canonical remainder modulo \(p\). For every \(f\in A[X]\),
\[
                   \ell_p(fV_w)=f(w).
\]
Repeated roots of \(p\) are allowed.
\end{proposition}

\begin{proof}
Synthetic division gives \((X-w)V_w=p\). The polynomial \(V_w\)
is monic of degree \(m-1\), so \(\ell_p(V_w)=1\).
For every polynomial \(g\), the identity
\((X-w)gV_w=pg\) and the vanishing of \(\ell_p\) on multiples
of \(p\) imply
\[
               \ell_p(XgV_w)=w\ell_p(gV_w).
\]
Induction gives \(\ell_p(X^jV_w)=w^j\). Linearity proves the
claim for arbitrary \(f\), with no restriction on its degree.
\end{proof}

The seven declarations in
\sourcefile{formal/VolumeLogic/ResolventFunctional.lean} prove this
identity for the constructed coefficient-list operations. For the
local matrix application, take the commutative algebra generated by
the one Casimir \(\Omega\) and scalar matrices. With
\(p=\chi_N\) and \(w=\Omega\), its \(V_w\) is exactly
\eqref{uf:eq-5}. This explains why the full functional recovers
\(f(\Omega)\), including its nilpotent part. The actual relation
\(\chi_N(\Omega)=0\) is supplied by
Proposition~\ref{prop:local-face-fourier}; the polynomial giving the
recorded crossing power is supplied by the separate crossing calculation.
Operators at different boxes
are not assumed to commute: their tensor contraction is multilinear
in the separately evaluated local matrices.

\section{Repeated roots, the odd simple root, and the midpoint
identity}\label{uf:sec-repeated-roots-the-odd-simple-root-and-the-midpoint-identity}

Write \begin{equation}\label{uf:eq-8}
c_a=2\cos\frac{(2a+1)\pi}{N},\qquad
 0\le a<\lfloor N/2\rfloor.
\end{equation} These are the distinct double roots. The only remaining root is
\(-2\), simple, when \(N\) is odd. At a double root \(c\), put
\(h(z)=\chi_N(z)/(z-c)^2\). The Chebyshev differential equation gives \begin{equation}\label{uf:eq-9}
h(c)=\frac{N^2}{4-c^2},\qquad
 \frac{h'(c)}{h(c)}=\frac{c}{4-c^2};
 \qquad \chi_N'(-2)=N^2\quad(N\ \mathrm{odd}).
\end{equation} Indeed \((4-z^2)\chi_N''-z\chi_N'+N^2(\chi_N-2)=0\). Evaluate it and
its derivative at \(c\), using \(\chi_N(c)=\chi_N'(c)=0\), to obtain
\eqref{uf:eq-9}.

If the resolvent at \(c\) is \(P_c/(z-c)+Q_c/(z-c)^2\), then
\(V_N(c)=h(c)Q_c\) and \(V_N'(c)=h(c)P_c+h'(c)Q_c\). Consequently the
\href{https://sashakolpakov.github.io/murakamis_dream/notes/actual_spectral_total_derivative_source.html}{original
spectral derivative functional}, on any scalar entry \(H\), is \begin{equation}\label{uf:eq-10}
\mathcal L_N^fH=
 \sum_a\frac{(4-c_a^2)(fH)'(c_a)-c_af(c_a)H(c_a)}{N^2}
 +\mathbf1_{N\ \mathrm{odd}}\frac{f(-2)H(-2)}{N^2}.
\end{equation} It sends \(V_N\) to \(f(\Omega)\), with \(f'(c)Q_c\) included. The
prime in \eqref{uf:eq-10} is a derivative in \(z\).

Let \begin{equation}\label{uf:eq-11}
\phi_k=\frac{(2k+1)\pi}{N},\quad 0\le k<N,\qquad
 P(\phi)=H(2\cos\phi).
\end{equation} Each double root appears twice in this grid. At the odd simple root,
\(\phi=\pi\) occurs once, and differentiating the sine supplies exactly
its weight in \eqref{uf:eq-10}. Thus \eqref{uf:eq-10} is equivalently \begin{equation}\label{uf:eq-12}
\mathcal L_N^fH
 =-\frac1{2}\sum_{k=0}^{N-1}
       \left.\frac{d}{d\phi}
       \frac{2\sin\phi\,f(2\cos\phi)H(2\cos\phi)}{N^2}
       \right|_{\phi=\phi_k}.
\end{equation} There is no omission or doubling of the odd simple node.

For the actual positive or negative integral crossing power, the
\href{https://sashakolpakov.github.io/murakamis_dream/notes/actual_direct_long_power_source.html}{full-Hermite
Gaussian} is \begin{equation}\label{uf:eq-13}
f_n(2\cos\phi)=q^{nd^2/4}
          \exp\left(\frac{inN(\phi-\pi)^2}{4\pi}\right),
 \qquad 0\le\phi\le2\pi.
\end{equation} Here \(q^a\) means \(e^{\pi ia/N}\), fixing the original square-root
lift even for nonintegral \(a\). In particular \begin{equation}\label{uf:eq-14}
|f_n|=1,\qquad |\partial_\phi f_n|\le\frac{|n|N}{2},\qquad
 f_n(c_a)=q^{nd^2/2}(-1)^{an}q^{na(a+1)}.
\end{equation} The nodal phase is the actual crossing phase, not its absolute value
substituted into the source. Formula \eqref{uf:eq-13} is symmetric about \(\pi\); it
has the required holomorphic germ at the odd simple node. The possible
endpoint cut at \(z=2\) does not occur at a spectral node and does not
enter the finite identity \eqref{uf:eq-12}.

Expanding \eqref{uf:eq-12}, with primes now in \(\phi\), gives \begin{equation}\label{uf:eq-15}
\mathcal L_N^{n}H=-\frac1{N^2}\sum_{k=0}^{N-1}
       \bigl(\cos\phi_k f_kP_k+
        \sin\phi_k f_k'P_k+\sin\phi_k f_kP_k'\bigr).
\end{equation} This finite formula is sufficient for every norm below; no
interchange of an infinite Fourier series or a contour limit is needed.

\section{An elementary exact derivative
inequality}\label{uf:sec-an-elementary-exact-derivative-inequality}

For every complex trigonometric polynomial \(P\) of degree at most
\(d\ge1\), set \(\theta_j=(2j-1)\pi/(2d)\), \(1\le j\le2d\). The finite
Riesz interpolation identity is \begin{equation}\label{uf:eq-16}
P'(x)=\sum_{j=1}^{2d}
       \frac{(-1)^{j+1}}{4d\sin^2(\theta_j/2)}P(x+\theta_j).
\end{equation} For completeness, invert the finite Fourier matrix on modes
\(-d+1,\ldots,d\). If \(z=e^{-i\theta_j}\), so \(z^{2d}=-1\),
differentiating the finite geometric sum gives its weight \begin{equation}\label{uf:eq-17}
\frac1{2d}\sum_{\ell=-d+1}^{d}i\ell e^{-i\ell\theta_j}
 =\frac{i z^{-d+1}}{d(1-z)^2}
 =\frac{(-1)^{j+1}}{4d\sin^2(\theta_j/2)}.
\end{equation} The extra mode \(-d\) has samples opposite to those of \(d\), and its
derivative at zero is opposite too. Thus \eqref{uf:eq-16} holds on all modes
\(-d,\ldots,d\); translation gives arbitrary \(x\).

The elementary sine product
\(\prod_{j=0}^{2d-1}\sin(x+j\pi/(2d))=C_d\sin(2dx)\), with \(C_d\ne0\),
yields after two logarithmic derivatives and substitution \(x=\pi/(4d)\)
\begin{equation}\label{uf:eq-18}
\sum_{j=1}^{2d}\frac1{4d\sin^2(\theta_j/2)}=d.
\end{equation} Taking absolute values in \eqref{uf:eq-16} therefore proves \begin{equation}\label{uf:eq-19}
\|P'\|_{\mathbb R}\le d\|P\|_{\mathbb R}.
\end{equation} The argument works without real coefficients or positivity.

The same finite identity and triangle inequality hold when the
coefficients lie in any finite-dimensional normed space. Thus \eqref{uf:eq-19}, and
consequently \eqref{uf:eq-1}--\eqref{uf:eq-4}, hold for such vector-valued polynomials without a
dimension factor. The intended application here remains the original
shared scalar coefficient \eqref{uf:eq-6}; no change of the source norm is required.

\section{Full and vanishing-node
bounds}\label{uf:sec-full-and-vanishing-node-bounds}

Since \(H\) has degree at most \(d\), \(P(\phi)=H(2\cos\phi)\) has
trigonometric degree at most \(d\), and its supremum is exactly
\(\|H\|_{[-2,2]}\). Equations \eqref{uf:eq-14}, \eqref{uf:eq-15}, and \eqref{uf:eq-19} give \begin{equation}\label{uf:eq-20}
|\mathcal L_N^{n}H|
 \le\frac{N}{N^2}\left(1+\frac{|n|N}{2}+d\right)\|H\|
 =\left(1+\frac{|n|}{2}\right)\|H\|,
\end{equation} because \(d+1=N\). This proves \eqref{uf:eq-1}.

If \(H\) vanishes at every node, the first two terms of \eqref{uf:eq-15} vanish
exactly. Only the derivative of the original numerator remains. Applying
\eqref{uf:eq-19} to that term gives \eqref{uf:eq-2}, uniformly in every signed \(n\), including
zero.

Here is why the original lower and upper one-sided numerators satisfy
this hypothesis. Put
\(D_-(z)=\prod_{a=0}^{\lfloor N/2\rfloor-1}(z-c_a)\). The exact
lower-row restriction leaves the factor \begin{equation}\label{uf:eq-21}
\frac{\chi_N(z)}{D_-(z)}=
 \begin{cases}D_-(z),&N\ \mathrm{even},\\
 (z+2)D_-(z),&N\ \mathrm{odd}.
 \end{cases}
\end{equation} Every lower one-sided entry is divisible by \eqref{uf:eq-21}; literal row
reflection gives the same statement for the upper piece. Thus values
vanish at all nodes. On the lower invariant block \(\Omega_-\), whose
eigenvalues \(c_a\) are distinct, \eqref{uf:eq-10} reduces exactly to \begin{equation}\label{uf:eq-22}
\mathcal L_N^nH_{V,\mathrm L}
 =\sum_a f_n(c_a)\frac{4-c_a^2}{N^2}
                       H_{V,\mathrm L}'(c_a).
\end{equation} These numerator derivatives are the original lower-block spectral
projectors. The derivative of \(f_n\) disappears, but the semisimple
source is not zero. For example the literal local diagonal coefficient
obeys \begin{equation}\label{uf:eq-23}
[c^d]H_{V,\mathrm L}(z)=\frac{\chi_N(z)}{z-c_0},
 \qquad [c^d]H_{R^n,\mathrm L}=q^{nd^2/2}.
\end{equation} Its modulus is one. This is an actual local coefficient, not a claim
that the same boundary can maximize a closed source.

\section{Shared polynomials and the unsplit complementary
source}\label{uf:sec-shared-polynomials-and-the-unsplit-complementary-source}

The other finite functionals preserve each partial node-vanishing
condition. Applying \eqref{uf:eq-1} and \eqref{uf:eq-2} successively to the original shared
polynomial proves \eqref{uf:eq-3}. Coefficient extraction in \eqref{uf:eq-6} is performed first
and is never replaced by independent local maximization.

For clarity, \eqref{uf:eq-4} needs only joint grid vanishing, not vanishing on every
node hyperplane. In each coordinate \eqref{uf:eq-15} is a sum of an evaluation term
and a term differentiating \(P\). Before summing its \(N\) grid entries,
the two coefficient bounds are \(1+|n_e|N/2\) and \(1\), respectively.
For a subset \(A\) of differentiated variables, iterating \eqref{uf:eq-19} gives \begin{equation}\label{uf:eq-24}
\|\partial_A\mathcal D(2\cos\boldsymbol\phi)\|
                  \le d^{|A|}\|\mathcal D\|_{[-2,2]^r}.
\end{equation} The term \(A=\varnothing\) vanishes pointwise at the tensor grid.
Every other term is bounded by
\(\prod_{e\in A}(d/N)\prod_{e\notin A}(1/N+|n_e|/2)\|\mathcal D\|\).
Adding the nonempty subsets proves \eqref{uf:eq-4}.

In the
\href{https://sashakolpakov.github.io/murakamis_dream/notes/actual_cross_cofactor_factorization.html}{original
central-cut decomposition}, the lower and upper pieces have the factor
\eqref{uf:eq-21}. The full numerator is the sum of those pieces and the central
piece defined by the endpoint conditions in Section~\ref{mca:sec-1}. At odd levels the
central row belongs to that crossing piece. Hence every local full
numerator agrees with its central piece at every spectral node. Multiplication and
the one original shared coefficient extraction show \begin{equation}\label{uf:eq-25}
(S_{\mathrm{full}}-S_{\mathrm C})
                    (c_{a_1},\ldots,c_{a_r})=0
\end{equation} for all choices of nodes, including odd simple nodes. Applying \eqref{uf:eq-4} to
this unsplit difference avoids summing the absolute values of
exponentially many path terms. The statement is finite and exact even
though its remaining shared supremum can be large.

\section{What the norm theorem does and does not
supply}\label{uf:sec-what-the-norm-theorem-does-and-does-not-supply}

There is already a level-independent dual filling cost in the
\href{https://sashakolpakov.github.io/murakamis_dream/notes/actual_filling_fourier_hermite_gram.html}{positive
Fourier--Hermite Gram construction}: its squared Hilbert dual norm is
exactly \(1+n^2(N^2-1)/(12N^2)\). That theorem uses a specified positive
energy of the entire mixed-color packet. The present bounds use instead
the supremum of the original shared resolvent numerator. The additional
statements here are the twist-independent vanishing-node estimate and
the unsplit product-grid cancellation. Neither construction by itself
supplies the geometric exponential bound for its remaining source norm.

For a fixed diagram and fixed powers, all constants in \eqref{uf:eq-1}--\eqref{uf:eq-4} are
independent of \(N\). Thus any sharp exponential estimate for the same
shared spectral supremum passes to its actual filled source without a
new exponential cost. The theorem also shows precisely where a
derivative of the crossing phase drops out and why this saves one or
more twist factors.

It does not prove that \(\|S_N\|\), or the complementary norm in \eqref{uf:eq-4},
grows at the geometric volume rate. Neither node vanishing nor the
uniform functional norm controls this remaining between-node polynomial.
Separate absolute local envelopes may exceed that rate, and \eqref{uf:eq-23}
prevents an argument that every one-sided channel is identically killed.
The theorem asserts no new saddle dominance, peripheral representation,
or knot-volume equality.

The \href{https://github.com/sashakolpakov/murakamis_dream/blob/main/experiments/audit_actual_uniform_filling_supremum_bound.py}{finite reproduction program}
takes \texttt{-\/-repo} pointing to the research checkout. It checks the
Chebyshev double/simple-root identities at levels 2--12 and the Riesz
formula at degrees 1--12 exactly in cyclotomic fields. It also compares
the midpoint reconstruction with independently built original
coproduct/Casimir and signed braid matrices at levels 3--8, for powers
\(\pm1,\ldots,\pm5\), and checks the lower-block compression including
its normal data. The JSON output distinguishes the exact rational
cyclotomic calculations from exact finite-field matrix checks. These
finite cases check algebraic identities; the all-level inequalities
follow from the proofs above, not from sampled norm ratios.

\section{The estimate needed for long
fillings}\label{uf:sec-the-estimate-needed-for-long-fillings}

It is useful to state the remaining comparison with its quantifiers. Fix
the diagram, and write \[
 \mathcal D_N=S_{N,\mathrm{full}}-S_{N,\mathrm C}.
\] This polynomial does not depend on the twist powers. Suppose one
proves, for some fixed \(\eta>0\), \begin{equation}\label{uf:eq-26}
\|\mathcal D_N\|_{[-2,2]^r}
       \le \exp\bigl(N(\Phi-\eta)+o(N)\bigr).
\end{equation} Here \(\Phi\) is the complete augmented action in the normalization
of the source. When it is identified geometrically,
\(\Phi=\operatorname{Vol}(M_{\mathrm{aug}})/(2\pi)\).

For each fixed tuple \(\boldsymbol n\), equation \eqref{uf:eq-4} then gives \begin{equation}\label{uf:eq-27}
\left|\left(\bigotimes_e\mathcal L_N^{n_e}\right)\mathcal D_N\right|
       \le \exp\bigl(N(\Phi-\eta)+o(N)\bigr).
\end{equation} The finite factor in \eqref{uf:eq-4} is absorbed in \(o(N)\); this argument does
not claim uniformity in twist powers growing with \(N\).

Suppose separately that the crossing contribution has a nonzero
asymptotic \begin{equation}\label{uf:eq-28}
Z_{N,\mathrm C}(\boldsymbol n)
       =a_N\exp(N\mathcal A_{\boldsymbol n}),\qquad
       N^{-1}\log|a_N|\longrightarrow0,
\end{equation} where \(\Re\mathcal A_{\boldsymbol n}>\Phi-\eta\). Equations \eqref{uf:eq-27} and
\eqref{uf:eq-28} show that the complementary contribution divided by the crossing
contribution tends to zero exponentially. Consequently \begin{equation}\label{uf:eq-29}
\lim_{N\to\infty}N^{-1}\log|Z_N(\boldsymbol n)|
       =\Re\mathcal A_{\boldsymbol n}.
\end{equation} The subexponential normalization factors in section 1 preserve this
limit. Identifying \(2\pi\Re\mathcal A_{\boldsymbol n}\) with the
hyperbolic volume still requires the geometric filling argument for that
diagram.

This implication explains why a strict gap in \eqref{uf:eq-26} is useful for
sufficiently long fixed fillings. The desired filled action approaches
the augmented action as the auxiliary fillings become long. A bound for
the full twist-independent polynomial at the augmented rate alone would
give only that larger rate; it would not prove the smaller filled
volume. The norm theorem proves the passage from \eqref{uf:eq-26} to \eqref{uf:eq-27}. It does
not prove \eqref{uf:eq-26} or supply \eqref{uf:eq-28} for a new diagram.
